% !TeX root = mainstarkov.tex
% !TeX program = xelatex
% !TeX encoding = UTF-8
\section{Анализ}

\subsection{Предметная область алгоритмов поиска путей}

Задача построения оптимального маршрута в замкнутом пространстве сводится к поиску кратчайшего пути на взвешенном неориентированном графе. Вершины графа соответствуют помещениям, коридорам, лестничным площадкам и лифтам, а рёбра -- проходам между ними с весами, пропорциональными расстоянию или времени прохождения. Такая модель является стандартной для систем indoor-навигации \cite{mongush2017} и опирается на классическую теорию графов \cite{zhang2015}.

В отличие от наружной навигации, где маршруты строятся по дорожной сети с относительно простой топологией, навигация внутри здания характеризуется многоэтажностью, наличием переходов между этажами с различным расстоянием (лестницы, лифты), а также возможностью динамического изменения проходимости участков. При возникновении чрезвычайной ситуации определённые проходы могут стать заблокированными, что требует немедленного перестроения маршрута. Существующие подходы к навигации внутри зданий охватывают как методы определения положения пользователя~-- по Bluetooth, Wi-Fi, компьютерному зрению, QR-меткам \cite{grinyak2018,singh2021,tomazic2021}, так и построение и сопровождение самого маршрута \cite{kulikov2021,landau2020,rubio2021,ohotichenko2021}; настоящая работа сосредоточена на втором. Особую роль навигация играет при эвакуации: требования к путям и выходам закреплены в нормативных документах \cite{fz123,sp113130}, информационные технологии всё активнее применяются для повышения пожарной безопасности зданий \cite{gusarova2023}, тогда как бумажные планы эвакуации статичны и требуют времени на изучение \cite{samoshin2026}.

Предметная область алгоритмов маршрутизации в системах эвакуации включает следующие ключевые задачи:
\begin{itemize}
\item поиск кратчайшего пути между двумя заданными точками внутри здания;
\item учёт весов рёбер, отражающих реальное расстояние или время прохождения;
\item динамическое исключение заблокированных узлов без модификации исходного графа;
\item поиск кратчайшего пути до ближайшего эвакуационного выхода из текущего положения;
\item обработка межэтажных переходов с учетом веса ребра;
\item обеспечение производительности, достаточной для работы в реальном времени.
\end{itemize}

\subsection{Обзор алгоритмов поиска кратчайших путей}

Для решения задачи поиска кратчайших путей на графах разработан ряд классических алгоритмов, каждый из которых имеет свои области применения, преимущества и ограничения. Ниже рассмотрены основные алгоритмы, применимые к задаче эвакуационной маршрутизации \cite{amelichev2022,isaev2023,levitin}.

\textbf{1. Алгоритм Дейкстры} \cite{dijkstra,mirahadi2021} находит кратчайшие пути от одной вершины до всех остальных в графе с неотрицательными весами рёбер. Алгоритм работает следующим образом: на каждом шаге выбирается непосещённая вершина с минимальным текущим расстоянием, после чего выполняется релаксация всех исходящих из неё рёбер. Сложность алгоритма составляет $O((V + E) \log V)$ при использовании бинарной кучи или $O(E + V \log V)$ с кучей Фибоначчи, где $V$ -- число вершин, $E$ -- число рёбер. Алгоритм гарантирует оптимальность решения для графов с неотрицательными весами и является одним из наиболее изученных и надёжных алгоритмов поиска путей.

Достоинства алгоритма Дейкстры: гарантия оптимальности при неотрицательных весах; эффективная работа на разреженных графах; возможность остановки при достижении целевой вершины; простота модификации для учёта динамических ограничений; широкая распространённость и наличие проверенных реализаций.

Недостатки: не работает с отрицательными весами рёбер; требует полной инициализации расстояний для всех вершин; не учитывает направление к цели (в отличие от эвристических алгоритмов).

\textbf{2. Алгоритм A*} является расширением алгоритма Дейкстры с использованием эвристической функции $h(v)$, оценивающей расстояние от вершины $v$ до цели. Функция приоритета имеет вид $f(v) = g(v) + h(v)$, где $g(v)$ -- текущее кратчайшее расстояние от старта, $h(v)$ -- эвристическая оценка расстояния до цели. При допустимой эвристике ($h(v) \leq$ реальное расстояние) A* гарантирует оптимальность и обычно посещает меньше вершин, чем Дейкстра.

Достоинства A*: значительное сокращение области поиска при подходящей эвристике; оптимальность при допустимой эвристике; широкое применение в игровой индустрии и робототехнике.

Недостатки: требует разработки эвристической функции, что затруднительно для indoor-графов с произвольной топологией; эвристика на основе евклидова расстояния может быть неточной из-за стен и непроходимых участков; при динамическом исключении узлов эвристика может стать недопустимой; производительность сильно зависит от качества эвристики.

\textbf{3. Алгоритм Флойда--Уоршелла} находит кратчайшие пути между всеми парами вершин за $O(V^3)$. Алгоритм последовательно рассматривает каждую вершину как промежуточную и обновляет расстояния между всеми парами. Результат хранится в матрице расстояний, что позволяет получать путь между любыми двумя вершинами за $O(1)$.

Достоинства: полная матрица кратчайших путей; простота реализации; работа с отрицательными весами (без отрицательных циклов).

Недостатки: кубическая сложность делает алгоритм неприменимым для графов с сотнями вершин; при изменении одного ребра требуется полный пересчёт; высокое потребление памяти ($O(V^2)$); не подходит для динамической маршрутизации.

\textbf{4. Поиск в ширину (BFS)} обходит граф уровнями от стартовой вершины, гарантируя нахождение кратчайшего пути в количестве рёбер. BFS работает за $O(V + E)$ и является оптимальным для невзвешенных графов.

Достоинства: линейная сложность; простота реализации; нахождение кратчайшего пути по числу рёбер.

Недостатки: не учитывает веса рёбер, что делает его непригодным для графов с различной длиной проходов; не различает короткие коридоры и длинные лестницы; не применим для оптимизации по расстоянию или времени.

\textbf{5. Алгоритм Беллмана--Форда} находит кратчайшие пути от одной вершины до всех остальных, допуская отрицательные веса рёбер. Сложность составляет $O(V \cdot E)$. Алгоритм выполняет $V - 1$ итераций релаксации всех рёбер.

Достоинства: поддержка отрицательных весов; обнаружение отрицательных циклов.

Недостатки: значительно медленнее Дейкстры ($O(V \cdot E)$ vs $O((V+E)\log V)$); отрицательные веса не требуются в задаче навигации; избыточность для рассматриваемой задачи.

\subsection{Сравнительный анализ алгоритмов}

Сравнение рассмотренных алгоритмов по ключевым критериям приведены в таблицах~\ref{algcompare:table1} и \ref{algcompare2:table1}.

\begin{xltabular}{\textwidth}{|X|c|c|c|}
\caption{Сравнение алгоритмов Дейкстры, A* и Флойда\label{algcompare:table1}}\\ \hline
\centrow Критерий & \centrow Дейкстра & \centrow A* & \centrow Флойд \\ \hline
\endfirsthead
\continuecaption{Продолжение таблицы \ref{algcompare:table1}}
\centrow Критерий & \centrow Дейкстра & \centrow A* & \centrow Флойд \\ \hline
\finishhead
Сложность & $O((V{+}E)\log V)$ & $O((V{+}E)\log V)$ & $O(V^3)$ \\ \hline
Оптимальность & Да & Да* & Да \\ \hline
Отриц.\ веса & Нет & Нет & Да \\ \hline
Динамич.\ искл.\ узлов & Да & Затруднено & Нет \\ \hline
Требует эвристику & Нет & Да & Нет \\ \hline
Память & $O(V{+}E)$ & $O(V{+}E)$ & $O(V^2)$ \\ \hline
Пересчёт при изменении & Полный & Полный & Полный
\end{xltabular}

\begin{xltabular}{\textwidth}{|X|c|c|}
\caption{Сравнение алгоритмов BFS и Беллмана--Форда\label{algcompare2:table1}}\\ \hline
\centrow Критерий & \centrow BFS & \centrow Беллман--Форд \\ \hline
\endfirsthead
\continuecaption{Продолжение таблицы \ref{algcompare2:table1}}
\centrow Критерий & \centrow BFS & \centrow Беллман--Форд \\ \hline
\finishhead
Сложность & $O(V{+}E)$ & $O(V \cdot E)$ \\ \hline
Оптимальность & По рёбрам & Да \\ \hline
Отриц.\ веса & Нет & Да \\ \hline
Динамич.\ искл.\ узлов & Да & Да \\ \hline
Требует эвристику & Нет & Нет \\ \hline
Память & $O(V{+}E)$ & $O(V{+}E)$ \\ \hline
Пересчёт при изменении & Полный & Полный
\end{xltabular}
\subsection{Представление графа в памяти}

Граф можно хранить двумя основными способами~-- матрицей смежности или списком смежности; выбор заметно влияет на расход памяти и скорость работы алгоритма.

\emph{Матрица смежности}~-- это таблица $V \times V$, где на пересечении строки $i$ и столбца $j$ стоит длина ребра между вершинами $i$ и $j$ (или признак его отсутствия). Проверка наличия ребра выполняется за $O(1)$, но память расходуется как $O(V^2)$ независимо от числа рёбер, а перебор соседей одной вершины стоит $O(V)$.

\emph{Список смежности} хранит для каждой вершины перечень её соседей с длинами. Память при этом~-- $O(V + E)$, а соседи перебираются ровно за их количество.

Граф здания разрежен: у каждого узла обычно два-четыре соседа, поэтому $E \approx 2{-}3V$. Сравнение двух представлений для такого случая приведено в таблице~\ref{graphrepr:table}.

\begin{xltabular}{\textwidth}{|X|c|c|}
	\caption{Сравнение способов представления графа\label{graphrepr:table}}\\ \hline
	\centrow Характеристика & \centrow Матрица & \centrow Список \\ \hline
	\endfirsthead
	\continuecaption{Продолжение таблицы \ref{graphrepr:table}}
	\centrow Характеристика & \centrow Матрица & \centrow Список \\ \hline
	\finishhead
	Память & $O(V^2)$ & $O(V{+}E)$ \\ \hline
	Перебор соседей вершины & $O(V)$ & $O(\deg v)$ \\ \hline
	Проверка наличия ребра & $O(1)$ & $O(\deg v)$ \\ \hline
	Подходит для разреженных графов & Нет & Да
\end{xltabular}

Алгоритм Дейкстры на каждом шаге перебирает соседей текущей вершины, а проверка отдельных рёбер ему почти не нужна. Поэтому для разреженного графа здания выбран список смежности: он экономит память (на 586 узлах матрица заняла бы около 340~тысяч ячеек против примерно 2,4~тысячи связей в списке) и ускоряет основной цикл алгоритма.

\subsection{Обоснование выбора алгоритма Дейкстры}

Для системы построения эвакуационных маршрутов выбран алгоритм Дейкстры по следующим причинам:

\begin{itemize}
\item эффективная работа на разреженных графах, характерных для планов зданий ($E \approx 2-3V$), в отличие от алгоритма Флойда--Уоршелла, требующего $O(V^3)$;
\item возможность модификации для исключения заблокированных узлов без полного пересчёта структуры графа -- достаточно передать множество исключений как параметр;
\item простота реализации с использованием стандартных структур данных .NET (\texttt{SortedSet}, \texttt{Dictionary}, \texttt{HashSet});
\item достаточная производительность для графов с сотнями узлов (целевое время $<200$ мс);
\item отсутствие необходимости в эвристической функции, разработка которой затруднена для indoor-графов с произвольной топологией;
\item возможность остановки при достижении целевой вершины без обработки всего графа.
\end{itemize}

Алгоритм A* не выбран, поскольку для навигационного графа здания с произвольной конфигурацией стен и переходов трудно построить допустимую эвристику. Евклидово расстояние не учитывает непроходимые стены, а манхэттенское расстояние не соответствует реальной топологии коридоров. Кроме того, при динамическом исключении узлов эвристика может стать недопустимой, что нарушит гарантию оптимальности.

\subsection{Особенности применения алгоритма Дейкстры к эвакуационной маршрутизации}

Применение алгоритма Дейкстры к задаче эвакуации требует учёта следующих специфических требований, отличающих данную задачу от стандартного поиска кратчайшего пути.

\textbf{Динамическое исключение узлов.} В режиме чрезвычайной ситуации или при обходе препятствий определённые узлы графа (проходы, двери, лестницы) могут стать недоступными. Алгоритм должен исключать такие узлы из рассмотрения без модификации исходного графа. Это реализуется путём передачи множества исключённых идентификаторов как параметра алгоритма. При релаксации рёбер алгоритм проверяет, не находится ли соседняя вершина в множестве исключений, и пропускает её. Конечная вершина маршрута никогда не исключается, даже если она присутствует в множестве \cite{mirahadi2020,zhangfire2019}.

\textbf{Вес межэтажных переходов.} Переход между этажами через лестницу или лифт имеет более большой вес ребра по сравнению с горизонтальным перемещением по коридору. Для межэтажных переходов назначается фиксированный вес, отражающий затраты времени на подъём/спуск. Данный подход позволяет алгоритму предпочитать горизонтальные маршруты, если они незначительно длиннее маршрута со сменой этажа \cite{gravit2019}.

\textbf{Поиск до ближайшего выхода.} Вместо поиска пути до заданной цели требуется найти кратчайший путь до любого из узлов, помеченных как эвакуационный выход. Это достигается запуском алгоритма Дейкстры от текущего положения пользователя с последующим сравнением расстояний до всех доступных выходов. Выбирается выход с минимальной суммарной стоимостью пути. Приоритет отдаётся выходам в том же здании, что и текущее положение пользователя \cite{han2020,yen2021}.

\textbf{Обработка отсутствия пути.} Если после завершения алгоритма расстояние до целевой вершины остаётся бесконечным (вершина недостижима), алгоритм должен возвращать пустой список, а не генерировать исключение. Это важно для устойчивости системы: при полной блокировке пути приложение должно корректно информировать пользователя, а не завершаться аварийно \cite{deng2021}.

\subsection{Анализ существующих реализаций}

На рынке существует ряд библиотек и фреймворков, предоставляющих реализацию алгоритмов поиска кратчайших путей. Готовые коммерческие платформы indoor-навигации \cite{navigineSdk,situmDocs,mappedinDocs,twogisFloors,mapstedOfficial} предоставляют навигацию внутри зданий как услугу, но они закрыты, требуют подписки и не дают доступа к реализации алгоритма маршрутизации. Поэтому ниже рассмотрены открытые алгоритмические библиотеки, наиболее известные из них.

\textbf{1. NetworkX} (Python) -- библиотека для работы с графами, предоставляющая реализации алгоритмов Дейкстры, A*, Флойда--Уоршелла и множества других. Широко используется в научных исследованиях и прототипировании. Однако библиотека написана на Python и не подходит для интеграции в .NET-приложение.

\textbf{2. QuickGraph} (C\#) -- библиотека графовых алгоритмов для .NET, предоставляющая реализации Дейкстры, A* и других алгоритмов. Поддерживает ориентированные и неориентированные графы с произвольными типами вершин. Недостатком является отсутствие поддержки динамического исключения узлов из коробки, а также отсутствие поддержки .NET MAUI.

\textbf{3. Boost.Graph} (C++) -- мощная библиотека графовых алгоритмов из набора Boost. Предоставляет высокопроизводительные реализации множества алгоритмов. Не подходит для .NET-приложения из-за языковых различий и сложности интеграции.

\textbf{4. Собственная реализация} -- разработка алгоритма Дейкстры на C\# с использованием стандартных коллекций .NET. Этот подход обеспечивает полный контроль над логикой алгоритма, возможность модификации для динамического исключения узлов, отсутствие внешних зависимостей и простоту отладки.

Для разрабатываемой системы выбрана собственная реализация, поскольку ни одна из существующих библиотек не предоставляет готовой поддержки динамического исключения узлов, необходимой для режима эвакуации. Кроме того, собственная реализация позволяет оптимизировать структуры данных под конкретную задачу и избежать лишних зависимостей.

\subsection{Выводы по разделу}

Проведённый анализ показывает, что алгоритм Дейкстры с использованием современных структур данных наилучшим образом соответствует требованиям задачи динамической эвакуационной маршрутизации. Его модификация для работы с динамически исключаемыми узлами и поиском ближайшего выхода позволяет создать надёжную основу для системы реального времени.

Ключевые результаты анализа:
\begin{itemize}
\item алгоритм Дейкстры обеспечивает оптимальность и достаточную производительность для графов зданий;
\item алгоритм A* не подходит из-за сложности построения допустимой эвристики для indoor-графов;
\item алгоритм Флойда--Уоршелла неприменим из-за кубической сложности;
\item BFS не учитывает веса рёбер;
\item собственная реализация предпочтительнее использования сторонних библиотек из-за необходимости поддержки динамического исключения узлов.
\end{itemize}
